\section{Current shared-split results and baseline audit}
\label{sec:shared-baselines}

This appendix reports the current, matched comparison across seven models.
It supersedes the older baseline comparisons that used an unprefixed neutral
prompt or the legacy split. All tables in this appendix are generated from
the retained measurements, rather than transcribed from historical reports.
Spearman correlation is the primary statistic; Pearson correlations are
included below. The historical main-text figures are not inputs to this analysis.

\subsection{Common evaluation protocol}

For each model and relation, we first remove every example whose object's
first answer token is shared by a \emph{different complete object} in that
relation. Multiple examples with the same complete object remain eligible.
This is a first-token identifiability filter, not a single-token-name filter:
multi-token subjects and multi-token objects can remain. After filtering,
seed 0 defines eight demonstrations, eight Jacobian references, and up to
100 test rows. The scored rows are disjoint from the reference rows.
Repeated subject strings can nevertheless occur in different rows; references
with the same subject string as a test row are excluded from that row's average.
We retain the existing 25 evaluable relations per model, including relations
with only two or three test rows, and separately report the sensitivity
requiring at least four. Every primary baseline receives the same test pool
and answer targets \emph{within each model}; undefined Geva scores are
handled explicitly below. Tokenizer-dependent object
filtering means that the examples need not be identical across models.

Every neutral lens uses exactly \texttt{\textbackslash n\{subject\} is known for},
with the tokenizer's native special-token behavior retained. The readout
position is the subject's final token. This avoids the Qwen3 8B first-position
artifact seen with the unprefixed prompt; the diagnostic did not find a
$+1$ or $-1$ indexing error. Qwen's original content tokens are unchanged by
the prefix. This is not universal: Mistral's initial word-boundary tokenization
changes, so its prefix comparison also changes tokenization. The
relation-specific eight-shot prompts, their model top-one answer tokens,
and the retained LRE measurements remain fixed. In particular, the newline
does not redefine the LRE target task.

All new lens measurements use float32, eager attention, frozen parameters,
and pinned model revisions. Hidden-state indices in
Table~\ref{tab:shared-coverage} denote the output after that many transformer
blocks. The layer is fixed across relations for a given model.
Jacobian products use double backward. Numerical checks compare final
logits, scalar and batched products, finite differences, and, where used,
cached-prefix derivatives. Each readout applies the model's actual final
normalization exactly once, followed by its output head and native logit
softcapping. This includes the Gemma/Qwen normalization parameterization and
replaces the historical OLMo raw-Jacobian file that predated the normalization fix.

For every test row, the target is the full model's top-one token under the
relation prompt, not necessarily the gold answer. LRE faithfulness is its
agreement with the selected relation-specific LRE. We retain the original
layer/beta selection and outcome values: this appendix does not turn that
in-sample hyperparameter selection into a held-out evaluation.

\begin{table}[htbp]
\centering
\caption{Coverage and the sensitivity excluding relations with fewer than four test rows. Empty Wikipedia retrievals remain in the primary evaluation. Undefined Geva rows have no predictions surviving the source filter and are excluded from attribute-rate means.}
\label{tab:shared-coverage}
\scriptsize
\setlength{\tabcolsep}{3.5pt}
\resizebox{\linewidth}{!}{%
\begin{tabular}{lrrrrrrr}
\toprule
Model & Layer & Rows & Subjects & Empty wiki & Undefined Geva & $n_{\geq4}$ & $\rho_{\geq4}$ \\
\midrule
OLMo 3 7B & 17 & 1544 & 1490 & 38 & 0 & 24 & +0.856 \\
GPT-J 6B & 15 & 1542 & 1479 & 48 & 0 & 24 & +0.749 \\
Mistral 7B & 17 & 1514 & 1459 & 42 & 0 & 24 & +0.628 \\
Gemma 2 9B & 22 & 1550 & 1488 & 43 & 0 & 24 & +0.591 \\
Qwen3 8B Base & 19 & 1544 & 1490 & 40 & 0 & 24 & +0.759 \\
Qwen3.8 27B & 34 & 1546 & 1488 & 48 & 0 & 24 & +0.644 \\
Pythia 6.9B & 17 & 1532 & 1473 & 42 & 413 & 24 & +0.631 \\
\bottomrule
\end{tabular}}
\end{table}


\subsection{Readout definitions}

Let $s_e$ be a test subject's intermediate state, $s_i$ a reference state,
$f_i(s_i)$ the last-position final-block state in reference $i$'s neutral
prompt, and $J_i=\partial f_i/\partial s_i$. Averages below use only the
eligible references for $e$. Let $D$ denote final normalization and vocabulary
readout. We compare
\begin{align*}
 \text{logit lens:}\quad &D(s_e),\\
 \text{differential logit lens:}\quad &D(s_e-\operatorname{mean}_i s_i),\\
 \text{raw Jacobian:}\quad &D(\operatorname{mean}_i J_i s_e),\\
 \text{differential Jacobian:}\quad &D(\operatorname{mean}_i J_i(s_e-s_i)),\\
 \text{affine Jacobian:}\quad
 &D\!\left(\operatorname{mean}_i[f_i(s_i)+\beta J_i(s_e-s_i)]\right).
\end{align*}
The raw Jacobian and affine readouts are distinct baselines. The fixed affine
baseline uses $\beta=1$; the full fixed grid is
$\{0.5,1,1.5,2,3,5\}$. We also retain both exploratory tuning rules: maximize
top-one agreement or minimize median log-rank within a relation, breaking
ties toward smaller beta. Both use the scored rows and are explicitly
in sample.

For neural readouts, a target's zero-based rank counts vocabulary logits
strictly greater than its logit. The relation score is
$-\operatorname{median}_e\log_{10}(\mathrm{rank}_e+1)$; higher is better.
Every relation receives equal weight in the correlation with faithfulness.
Ranks are not pooled over subjects or over models.

\subsection{Wikipedia attribute and frequency baselines}

The Geva baseline is the subject-attribute overlap metric of
\citet{geva2023dissecting}, using the retrieval/filtering reconstruction of
\citet{gerstner2026weakening}. It is not a frequency-rank correlation.
We use the verified English Wikipedia dump of 20 October 2021, containing
40,284,960 cleaned nonempty paragraphs. We rebuild the BM25 index for the
union of the current test subjects across all seven models. This common
query vocabulary is an explicit adaptation of the reconstruction: retaining
the old index would retain its obsolete subject vocabulary and document
lengths. Lucene BM25 uses $k_1=1.5$, $b=0.75$, and retrieves 100 paragraphs
before applying the source subject/title substring filter.

Each model tokenizes the concatenated retained paragraphs with its own
tokenizer. The attribute set keeps decoded tokens of stripped length greater
than two, excluding English stopwords. We take the \emph{top 50} ordinary
logit-lens predictions first and then apply the source prediction filter.
The subject's attribute rate is the overlap with the attribute set divided
by the number of surviving predictions, times 100. The source's asymmetric
case handling is preserved: the corpus-side stopword check is case-sensitive;
the prediction-side check lowercases. Relation scores average subject rates.
An empty retrieval contributes zero when predictions survive; an empty
prediction denominator is undefined, not silently set to zero. Coverage is
recorded separately. Pythia has 413 such undefined rows, which are excluded
from its attribute-rate means; the other models have none.
Table~\ref{tab:shared-geva-coverage} additionally restricts all readouts and
the LRE outcome to the same Geva-defined rows, with at least two per relation.
This transfers the metric to our models, layer, prompts,
and relation pool; it does not reproduce the source papers' original experiments.

The separate \emph{Wikipedia token-frequency rank} baseline ranks each
relation-prompt target by its count in the same subject's retained Wikipedia
text. Ties receive zero-based midranks across the full model output
vocabulary; absent tokens tie at the bottom. We aggregate with the same
negative median log-rank. Finally, the vocabulary-agreement diagnostic
computes Spearman correlation between Wikipedia token counts and logit-lens
logits over the eligible vocabulary, then averages over a relation's test
rows. This diagnostic is distinct from both Geva overlap and answer-token rank.

\subsection{Primary Spearman comparisons}

The differential-Jacobian correlation ranges from 0.479 to 0.729. The raw Jacobian has a higher point estimate in 5 of seven models; thus these results do not establish a general benefit from reference subtraction. For Qwen3 8B Base, ordinary logit lens reaches $\rho=0.743$, compared with $\rho=0.729$ for differential Jacobian. Every paired 95\% bootstrap interval for differential minus raw Jacobian includes zero. The tables retain negative as well as positive correlations and do not select a different baseline for each model after observing the results.

\begin{table}[htbp]
\centering
\caption{Spearman $\rho$ with shared-split LRE faithfulness across 25 relations per model. Each relation has equal weight; larger lens scores mean better target ranks.}
\label{tab:shared-lenses-spearman}
\scriptsize
\setlength{\tabcolsep}{3.5pt}
\resizebox{\linewidth}{!}{%
\begin{tabular}{lrrrrrrr}
\toprule
Predictor & OLMo & GPT-J & Mistral & Gemma & Qwen 8B & Qwen 27B & Pythia \\
\midrule
Differential Jacobian & +0.683 & +0.573 & +0.606 & +0.554 & +0.729 & +0.600 & +0.479 \\
Raw Jacobian & +0.689 & +0.567 & +0.586 & +0.575 & +0.736 & +0.644 & +0.501 \\
Logit lens & +0.234 & +0.531 & +0.402 & +0.554 & +0.743 & +0.402 & -0.343 \\
Differential logit lens & +0.326 & +0.394 & +0.249 & +0.221 & +0.607 & +0.454 & +0.217 \\
Affine Jacobian (beta=1) & +0.562 & +0.430 & +0.407 & +0.472 & +0.300 & +0.308 & +0.308 \\
Geva attribute rate & -0.114 & +0.002 & -0.304 & -0.292 & +0.052 & -0.196 & +0.096 \\
Wikipedia token-frequency rank & +0.277 & +0.140 & +0.276 & +0.419 & +0.374 & +0.304 & +0.154 \\
Wikipedia TF/logit vocabulary agreement & +0.301 & -0.109 & -0.168 & -0.102 & -0.175 & +0.065 & +0.310 \\
\bottomrule
\end{tabular}}
\end{table}

\begin{table}[htbp]
\centering
\caption{Two-sided, unadjusted Spearman $p$-values for the lenses comparisons (25 relations). These are exploratory tests, not multiplicity-adjusted claims.}
\label{tab:shared-lenses-pvalues}
\scriptsize
\setlength{\tabcolsep}{3.5pt}
\resizebox{\linewidth}{!}{%
\begin{tabular}{lrrrrrrr}
\toprule
Predictor & OLMo & GPT-J & Mistral & Gemma & Qwen 8B & Qwen 27B & Pythia \\
\midrule
Differential Jacobian & 0.00017 & 0.00273 & 0.00131 & 0.00409 & 3.56e-05 & 0.00152 & 0.0153 \\
Raw Jacobian & 0.000138 & 0.00315 & 0.0021 & 0.00262 & 2.77e-05 & 0.00051 & 0.0107 \\
Logit lens & 0.26 & 0.00634 & 0.0466 & 0.0041 & 2.13e-05 & 0.0461 & 0.0937 \\
Differential logit lens & 0.112 & 0.0514 & 0.231 & 0.289 & 0.00131 & 0.0225 & 0.297 \\
Affine Jacobian (beta=1) & 0.00344 & 0.032 & 0.0435 & 0.0173 & 0.145 & 0.135 & 0.135 \\
Geva attribute rate & 0.588 & 0.992 & 0.14 & 0.157 & 0.806 & 0.347 & 0.647 \\
Wikipedia token-frequency rank & 0.18 & 0.503 & 0.182 & 0.037 & 0.0652 & 0.14 & 0.463 \\
Wikipedia TF/logit vocabulary agreement & 0.144 & 0.603 & 0.422 & 0.627 & 0.403 & 0.759 & 0.132 \\
\bottomrule
\end{tabular}}
\end{table}

\begin{table}[htbp]
\centering
\caption{Geva coverage sensitivity: restrict every readout and the LRE outcome to the same rows with a defined Geva rate, retaining relations with at least two such rows. Entries are Spearman correlations. Only Pythia loses rows to the Geva prediction filter.}
\label{tab:shared-geva-coverage}
\scriptsize
\setlength{\tabcolsep}{3.5pt}
\resizebox{\linewidth}{!}{%
\begin{tabular}{lrrrrrrr}
\toprule
Model & Relations & Rows & Differential & Raw J & Logit lens & Geva & Wiki TF \\
\midrule
OLMo 3 7B & 25 & 1544 & +0.683 & +0.689 & +0.234 & -0.114 & +0.277 \\
GPT-J 6B & 25 & 1542 & +0.573 & +0.567 & +0.531 & +0.002 & +0.140 \\
Mistral 7B & 25 & 1514 & +0.606 & +0.586 & +0.402 & -0.304 & +0.276 \\
Gemma 2 9B & 25 & 1550 & +0.554 & +0.575 & +0.554 & -0.292 & +0.419 \\
Qwen3 8B Base & 25 & 1544 & +0.729 & +0.736 & +0.743 & +0.052 & +0.374 \\
Qwen3.8 27B & 25 & 1546 & +0.600 & +0.644 & +0.402 & -0.196 & +0.304 \\
Pythia 6.9B & 25 & 1119 & +0.393 & +0.394 & -0.116 & +0.204 & +0.029 \\
\bottomrule
\end{tabular}}
\end{table}

\begin{table}[htbp]
\centering
\caption{Spearman $\rho$ with shared-split LRE faithfulness across 25 relations per model. Tuned rows select beta on the scored test rows and are exploratory.}
\label{tab:shared-beta-spearman}
\scriptsize
\setlength{\tabcolsep}{3.5pt}
\resizebox{\linewidth}{!}{%
\begin{tabular}{lrrrrrrr}
\toprule
Predictor & OLMo & GPT-J & Mistral & Gemma & Qwen 8B & Qwen 27B & Pythia \\
\midrule
Affine beta=0.5 & +0.367 & +0.160 & +0.382 & +0.498 & +0.157 & +0.289 & +0.103 \\
Affine beta=1 & +0.562 & +0.430 & +0.407 & +0.472 & +0.300 & +0.308 & +0.308 \\
Affine beta=1.5 & +0.635 & +0.473 & +0.425 & +0.528 & +0.396 & +0.326 & +0.391 \\
Affine beta=2 & +0.658 & +0.486 & +0.420 & +0.585 & +0.507 & +0.364 & +0.428 \\
Affine beta=3 & +0.728 & +0.463 & +0.399 & +0.647 & +0.600 & +0.393 & +0.510 \\
Affine beta=5 & +0.746 & +0.481 & +0.459 & +0.657 & +0.675 & +0.447 & +0.537 \\
Affine tuned top1 (in sample) & +0.364 & +0.175 & +0.383 & +0.498 & +0.210 & +0.279 & +0.161 \\
Affine tuned median-logrank (in sample) & +0.746 & +0.475 & +0.459 & +0.657 & +0.666 & +0.427 & +0.510 \\
\bottomrule
\end{tabular}}
\end{table}

\begin{table}[htbp]
\centering
\caption{Paired differences in Spearman correlation: differential Jacobian minus each baseline, with 95\% relation-bootstrap intervals. Positive values favor differential Jacobian. All six baseline contrasts are retained in the accompanying CSV.}
\label{tab:shared-paired-gaps}
\scriptsize
\setlength{\tabcolsep}{3.5pt}
\resizebox{\linewidth}{!}{%
\begin{tabular}{lrrr}
\toprule
Model & Minus raw Jacobian & Minus logit lens & Minus affine \\
\midrule
OLMo 3 7B & -0.007 [-0.064, +0.039] & +0.449 [+0.071, +0.870] & +0.120 [-0.345, +0.527] \\
GPT-J 6B & +0.007 [-0.040, +0.055] & +0.043 [-0.237, +0.334] & +0.144 [-0.285, +0.565] \\
Mistral 7B & +0.021 [-0.132, +0.181] & +0.205 [-0.042, +0.482] & +0.199 [-0.283, +0.683] \\
Gemma 2 9B & -0.022 [-0.117, +0.074] & +0.000 [-0.299, +0.291] & +0.082 [-0.383, +0.526] \\
Qwen3 8B Base & -0.007 [-0.059, +0.032] & -0.014 [-0.180, +0.140] & +0.429 [+0.041, +0.812] \\
Qwen3.8 27B & -0.044 [-0.145, +0.035] & +0.198 [-0.130, +0.567] & +0.292 [-0.191, +0.772] \\
Pythia 6.9B & -0.022 [-0.149, +0.081] & +0.822 [+0.117, +1.349] & +0.172 [-0.238, +0.554] \\
\bottomrule
\end{tabular}}
\end{table}


The $p$-values are two-sided and unadjusted. Relations can share subjects
and objects, and many predictors are tested; these are descriptive,
exploratory comparisons. Paired bootstrap intervals resample the same 25
relations for both readouts, rerank each resample, and use 10,000 draws
with seed 20260921. They quantify relation-resampling uncertainty, not
variation across model training runs.

\subsection{Frequency, cardinality, and output-distribution controls}

The model answer prior is the mean log probability of a target token under
ten fixed neutral contexts: empty/BOS, ``The'', ``In the'', ``It was'',
``She said the'', ``Here is'', ``One of the'', ``According to the'',
``There is a'', and ``Yesterday the''. We average these token scores over
the current test rows, separately for model targets and gold first tokens.
These contexts define a frequency proxy and are not the neutral lens prompt.
They were rerun in float32 for every pinned model. The empty context uses
the tokenizer's BOS string, or its existing end-of-text fallback, and retains
the tokenizer's normal special-token insertion.

Cardinality is samples per distinct complete object. We report the original
full-data version as well as filtered-pool and test-only versions. The
relation-prompt intercept-only baseline decodes the mean reference output
without a subject-dependent Jacobian term. Its saved ordinary-forward
measurements are reused only after checking model revision, split fingerprint,
ordered target IDs, gold IDs, and selected-LRE agreement against the current
rows. Modal-target fraction, target entropy, and model gold-token accuracy
are outcome-distribution controls. Modal-target fraction is an in-sample
constant-prediction oracle, not a learned predictor.

\begin{table}[htbp]
\centering
\caption{Spearman $\rho$ with shared-split LRE faithfulness across 25 relations per model. Each relation has equal weight; larger lens scores mean better target ranks.}
\label{tab:shared-controls-spearman}
\scriptsize
\setlength{\tabcolsep}{3.5pt}
\resizebox{\linewidth}{!}{%
\begin{tabular}{lrrrrrrr}
\toprule
Predictor & OLMo & GPT-J & Mistral & Gemma & Qwen 8B & Qwen 27B & Pythia \\
\midrule
Model answer-token prior & -0.137 & +0.067 & -0.080 & +0.283 & +0.001 & +0.171 & -0.223 \\
Gold answer-token prior & +0.027 & +0.175 & +0.056 & +0.358 & +0.167 & +0.252 & +0.010 \\
Full-data samples/objects & +0.304 & +0.258 & +0.743 & +0.666 & +0.362 & +0.498 & +0.249 \\
Filtered samples/objects & +0.303 & +0.265 & +0.743 & +0.664 & +0.361 & +0.499 & +0.255 \\
Test samples/objects & +0.291 & +0.126 & +0.699 & +0.665 & +0.338 & +0.498 & +0.198 \\
Relation-prompt intercept-only agreement & +0.293 & +0.172 & +0.682 & +0.600 & +0.346 & +0.530 & +0.283 \\
Modal-target fraction (oracle) & +0.584 & +0.494 & +0.804 & +0.768 & +0.549 & +0.709 & +0.413 \\
Target entropy & -0.753 & -0.699 & -0.777 & -0.749 & -0.736 & -0.762 & -0.764 \\
Model gold-token accuracy & +0.683 & +0.781 & +0.335 & +0.370 & +0.549 & +0.447 & +0.775 \\
\bottomrule
\end{tabular}}
\end{table}

\begin{table}[htbp]
\centering
\caption{Two-sided, unadjusted Spearman $p$-values for the controls comparisons (25 relations). These are exploratory tests, not multiplicity-adjusted claims.}
\label{tab:shared-controls-pvalues}
\scriptsize
\setlength{\tabcolsep}{3.5pt}
\resizebox{\linewidth}{!}{%
\begin{tabular}{lrrrrrrr}
\toprule
Predictor & OLMo & GPT-J & Mistral & Gemma & Qwen 8B & Qwen 27B & Pythia \\
\midrule
Model answer-token prior & 0.515 & 0.749 & 0.705 & 0.17 & 0.997 & 0.415 & 0.284 \\
Gold answer-token prior & 0.898 & 0.402 & 0.791 & 0.0793 & 0.424 & 0.224 & 0.964 \\
Full-data samples/objects & 0.139 & 0.212 & 2.14e-05 & 0.000282 & 0.0755 & 0.0114 & 0.23 \\
Filtered samples/objects & 0.141 & 0.2 & 2.14e-05 & 0.000295 & 0.0766 & 0.0111 & 0.218 \\
Test samples/objects & 0.158 & 0.548 & 0.000101 & 0.000285 & 0.0984 & 0.0113 & 0.343 \\
Relation-prompt intercept-only agreement & 0.155 & 0.411 & 0.000172 & 0.00152 & 0.0901 & 0.00639 & 0.171 \\
Modal-target fraction (oracle) & 0.0022 & 0.0121 & 1.32e-06 & 7.33e-06 & 0.0045 & 7.22e-05 & 0.0404 \\
Target entropy & 1.43e-05 & 0.000103 & 4.83e-06 & 1.64e-05 & 2.75e-05 & 9.45e-06 & 8.69e-06 \\
Model gold-token accuracy & 0.000169 & 4.15e-06 & 0.101 & 0.0689 & 0.00452 & 0.0252 & 5.28e-06 \\
\bottomrule
\end{tabular}}
\end{table}


The corpus controls reuse a completed scan of the three published OLMo 3
training mixes: 175,877 shards and 3,861,363,584 document records. Of these,
2,329,747 records have redacted text that cannot contribute observed mentions.
These are OLMo training-corpus measurements; for the six other models, the
\emph{same OLMo corpus is only an external frequency proxy}. We do not claim
access to those models' actual training frequencies. Counts describe the
published document multiset, including repeated records across stages,
rather than packed training sequences or exact training exposure weights.

Counts are binary document presence of literal subject/object strings,
using NFC normalization, whitespace collapse, casefolding, and Unicode word
boundaries, without alias resolution or a proximity constraint. With joint
count $c_{so}$ and marginals $c_s,c_o$, the directional quantities are
$P(o|s)=c_{so}/c_s$, $P(s|o)=c_{so}/c_o$, and their difference. Both directions
use common support where both marginals are at least one. Macro features
average fact-level values on the exact current test rows; micro conditionals
are ratios of summed counts. Repeated evaluation rows retain their evaluation
weight even when they share a deduplicated corpus fact. The count scan itself
does not depend on the lens prompt and is reused; all subset aggregation and
correlations are recomputed. The overlap sensitivity removes pairs where one
normalized name contains the other. Conditional direction is not independent
of marginal frequency: when $c_{so}>0$, $P(o|s)/P(s|o)=c_o/c_s$.

\begin{table}[htbp]
\centering
\caption{Spearman $\rho$ with shared-split LRE faithfulness across 25 relations per model. Document counts come from OLMo's published training split; they are an external corpus proxy for the other models.}
\label{tab:shared-corpus-spearman}
\scriptsize
\setlength{\tabcolsep}{3.5pt}
\resizebox{\linewidth}{!}{%
\begin{tabular}{lrrrrrrr}
\toprule
Predictor & OLMo & GPT-J & Mistral & Gemma & Qwen 8B & Qwen 27B & Pythia \\
\midrule
Subject document frequency & +0.290 & +0.336 & -0.166 & -0.204 & +0.166 & +0.061 & +0.378 \\
Object document frequency & +0.644 & +0.642 & +0.618 & +0.590 & +0.603 & +0.627 & +0.610 \\
Joint document frequency & +0.578 & +0.620 & +0.158 & +0.232 & +0.489 & +0.382 & +0.631 \\
Mean log(1+joint frequency) & +0.568 & +0.590 & +0.147 & +0.239 & +0.478 & +0.377 & +0.607 \\
Log(1+mean joint frequency) & +0.578 & +0.620 & +0.158 & +0.232 & +0.489 & +0.382 & +0.631 \\
P(object | subject) & +0.386 & +0.418 & +0.566 & +0.635 & +0.411 & +0.418 & +0.305 \\
P(subject | object) & -0.328 & -0.287 & -0.740 & -0.654 & -0.421 & -0.502 & -0.230 \\
Forward minus reverse conditional & +0.362 & +0.415 & +0.636 & +0.687 & +0.432 & +0.452 & +0.288 \\
Micro P(object | subject) & +0.615 & +0.574 & +0.639 & +0.708 & +0.597 & +0.484 & +0.513 \\
Micro P(subject | object) & +0.030 & +0.014 & -0.485 & -0.361 & -0.095 & -0.213 & +0.021 \\
\bottomrule
\end{tabular}}
\end{table}

\begin{table}[htbp]
\centering
\caption{Two-sided, unadjusted Spearman $p$-values for the corpus comparisons (25 relations). These are exploratory tests, not multiplicity-adjusted claims.}
\label{tab:shared-corpus-pvalues}
\scriptsize
\setlength{\tabcolsep}{3.5pt}
\resizebox{\linewidth}{!}{%
\begin{tabular}{lrrrrrrr}
\toprule
Predictor & OLMo & GPT-J & Mistral & Gemma & Qwen 8B & Qwen 27B & Pythia \\
\midrule
Subject document frequency & 0.16 & 0.101 & 0.428 & 0.328 & 0.428 & 0.773 & 0.0621 \\
Object document frequency & 0.000508 & 0.000539 & 0.000987 & 0.0019 & 0.00142 & 0.000791 & 0.0012 \\
Joint document frequency & 0.00247 & 0.000943 & 0.45 & 0.264 & 0.0131 & 0.0595 & 0.000716 \\
Mean log(1+joint frequency) & 0.00305 & 0.0019 & 0.484 & 0.249 & 0.0157 & 0.0632 & 0.00131 \\
Log(1+mean joint frequency) & 0.00247 & 0.000943 & 0.45 & 0.264 & 0.0131 & 0.0595 & 0.000716 \\
P(object | subject) & 0.0564 & 0.0375 & 0.00322 & 0.000654 & 0.041 & 0.0377 & 0.138 \\
P(subject | object) & 0.109 & 0.164 & 2.33e-05 & 0.000392 & 0.0359 & 0.0105 & 0.269 \\
Forward minus reverse conditional & 0.0753 & 0.0389 & 0.000627 & 0.000147 & 0.0309 & 0.0233 & 0.162 \\
Micro P(object | subject) & 0.00107 & 0.00272 & 0.000584 & 7.45e-05 & 0.00162 & 0.0143 & 0.00874 \\
Micro P(subject | object) & 0.887 & 0.946 & 0.014 & 0.0762 & 0.651 & 0.308 & 0.921 \\
\bottomrule
\end{tabular}}
\end{table}

\begin{table}[htbp]
\centering
\caption{OLMo-only Spearman sensitivities by corpus stage and literal-overlap exclusion; relation counts in parentheses. All features share support where both marginals are observed.}
\label{tab:shared-corpus-stages}
\scriptsize
\setlength{\tabcolsep}{3.5pt}
\resizebox{\linewidth}{!}{%
\begin{tabular}{lrrrrr}
\toprule
Corpus subset & Log joint & Object freq. & $P(o|s)$ & $P(s|o)$ & Difference \\
\midrule
All stages & +0.568 (25) & +0.644 (25) & +0.386 (25) & -0.328 (25) & +0.362 (25) \\
All stages (no overlap) & +0.568 (25) & +0.645 (25) & +0.408 (25) & -0.328 (25) & +0.378 (25) \\
dolmino\_mix-100B & +0.659 (25) & +0.667 (25) & +0.371 (25) & -0.349 (25) & +0.380 (25) \\
dolmino\_mix-100B (no overlap) & +0.663 (25) & +0.670 (25) & +0.384 (25) & -0.362 (25) & +0.385 (25) \\
longmino\_mix-50B & +0.626 (25) & +0.680 (25) & +0.415 (25) & -0.264 (25) & +0.387 (25) \\
longmino\_mix-50B (no overlap) & +0.638 (25) & +0.676 (25) & +0.421 (25) & -0.288 (25) & +0.396 (25) \\
mix-6T & +0.568 (25) & +0.644 (25) & +0.386 (25) & -0.324 (25) & +0.362 (25) \\
mix-6T (no overlap) & +0.568 (25) & +0.645 (25) & +0.408 (25) & -0.328 (25) & +0.378 (25) \\
\bottomrule
\end{tabular}}
\end{table}


\subsection{Matched prominence and partial-correlation diagnostics}

We refresh the original OLMo/GPT-J prominence diagnostic on its 14-relation
scope. For up to the first 40 test rows, the model greedily generates up to
90 tokens from each of two newline-prefixed prose prompts:
\texttt{\textbackslash n\{subject\} is best known for} and
\texttt{\textbackslash nHere are some facts about \{subject\}. First,}.
A row is a hit when either continuation contains a distinctive gold-object
word under the existing matching rule. Words shared with the subject name
are removed; rows with no distinctive word are excluded. A separate
four-character-prefix rule flags residual name leakage. We require at least
five scorable rows per relation/subset. The gold-in-lens top-five rate and
selected-LRE faithfulness are recomputed on exactly these same rows, fixing
the old comparison's cohort mismatch. This diagnostic uses two different
prose prompts by design and is separate from the 25-relation lens benchmark.

\begin{table}[htbp]
\centering
\caption{Fresh free-generation prominence diagnostic on identical generation/lens rows. Entries show Spearman $\rho$ ($p$); LRE faithfulness is recomputed on those rows with the original selected LRE settings.}
\label{tab:shared-prominence}
\scriptsize
\setlength{\tabcolsep}{3.5pt}
\resizebox{\linewidth}{!}{%
\begin{tabular}{lrrrr}
\toprule
Model & Subset & Relations & vs. LRE faithfulness & vs. gold in lens top 5 \\
\midrule
OLMo 3 7B & All & 14 & +0.723 (0.00349) & +0.615 (0.0193) \\
OLMo 3 7B & No name leakage & 13 & +0.616 (0.025) & +0.610 (0.0269) \\
GPT-J 6B & All & 14 & +0.082 (0.781) & +0.494 (0.0726) \\
GPT-J 6B & No name leakage & 13 & -0.057 (0.854) & +0.425 (0.148) \\
\bottomrule
\end{tabular}}
\end{table}

\begin{table}[htbp]
\centering
\caption{Sensitivity of the differential-Jacobian association. Partial Spearman residualizes ranked variables on ranked controls; the last row changes the outcome to intercept-adjusted agreement.}
\label{tab:shared-sensitivity}
\scriptsize
\setlength{\tabcolsep}{3.5pt}
\resizebox{\linewidth}{!}{%
\begin{tabular}{lrrrrrrr}
\toprule
Control / outcome & OLMo & GPT-J & Mistral & Gemma & Qwen 8B & Qwen 27B & Pythia \\
\midrule
cardinality & +0.718 & +0.610 & +0.582 & +0.611 & +0.754 & +0.684 & +0.547 \\
answer\_prior & +0.675 & +0.579 & +0.604 & +0.505 & +0.730 & +0.598 & +0.451 \\
cardinality + answer\_prior & +0.709 & +0.619 & +0.578 & +0.591 & +0.755 & +0.682 & +0.527 \\
object\_frequency & +0.605 & +0.482 & +0.367 & +0.406 & +0.654 & +0.501 & +0.433 \\
cardinality + answer\_prior + object\_frequency & +0.540 & +0.403 & +0.379 & +0.546 & +0.639 & +0.583 & +0.315 \\
Intercept-adjusted outcome & +0.713 & +0.665 & +0.683 & +0.641 & +0.831 & +0.721 & +0.630 \\
\bottomrule
\end{tabular}}
\end{table}

\begin{table}[htbp]
\centering
\caption{Differential-Jacobian robustness: leave-one-domain-out Spearman range and median, two-sided permutation $p$ (10,000 draws, plus-one correction), and correlation after restricting the LRE outcome to at least four model-correct test rows per relation. The last column retains the original full-test lens score; relation counts are in parentheses.}
\label{tab:shared-robustness}
\scriptsize
\setlength{\tabcolsep}{3.5pt}
\resizebox{\linewidth}{!}{%
\begin{tabular}{lrrrr}
\toprule
Model & Domain range & Median & Permutation $p$ & Correct-only outcome \\
\midrule
OLMo 3 7B & 0.589--0.855 & 0.675 & 0.0002 & +0.841 (24) \\
GPT-J 6B & 0.421--0.721 & 0.569 & 0.0032 & +0.794 (24) \\
Mistral 7B & 0.550--0.648 & 0.604 & 0.0015 & +0.655 (24) \\
Gemma 2 9B & 0.363--0.721 & 0.553 & 0.0047 & +0.612 (24) \\
Qwen3 8B Base & 0.658--0.808 & 0.720 & 9.999e-05 & +0.726 (24) \\
Qwen3.8 27B & 0.556--0.658 & 0.589 & 0.0018 & +0.580 (24) \\
Pythia 6.9B & 0.286--0.577 & 0.472 & 0.0162 & +0.516 (24) \\
\bottomrule
\end{tabular}}
\end{table}


Partial Spearman correlation here is Pearson correlation between residuals
after regressing ranked lens score and ranked faithfulness on an intercept
and the ranked controls. It is not Spearman correlation of reranked
residuals. The intercept-adjusted outcome is
$(\mathrm{faith}-\mathrm{prior\ faith})/(1-\mathrm{prior\ faith})$;
relations with zero denominator are excluded and their counts are retained
in the machine-readable report.

\clearpage
\subsection{Pearson sensitivity}
These tables change only the correlation statistic. The samples, feature
definitions, transformations, and LRE outcomes are identical to the
Spearman tables. Pearson is computed on the already defined lens scores
and aggregate corpus features, not on raw subject ranks.
\begin{table}[htbp]
\centering
\caption{Pearson $r$ with shared-split LRE faithfulness across 25 relations per model. Each relation has equal weight; larger lens scores mean better target ranks.}
\label{tab:shared-lenses-pearson}
\scriptsize
\setlength{\tabcolsep}{3.5pt}
\resizebox{\linewidth}{!}{%
\begin{tabular}{lrrrrrrr}
\toprule
Predictor & OLMo & GPT-J & Mistral & Gemma & Qwen 8B & Qwen 27B & Pythia \\
\midrule
Differential Jacobian & +0.737 & +0.533 & +0.627 & +0.509 & +0.740 & +0.536 & +0.582 \\
Raw Jacobian & +0.728 & +0.528 & +0.648 & +0.561 & +0.759 & +0.576 & +0.599 \\
Logit lens & +0.247 & +0.507 & +0.510 & +0.559 & +0.676 & +0.398 & -0.237 \\
Differential logit lens & +0.333 & +0.493 & +0.341 & +0.189 & +0.551 & +0.331 & +0.375 \\
Affine Jacobian (beta=1) & +0.544 & +0.275 & +0.425 & +0.510 & +0.408 & +0.370 & +0.298 \\
Geva attribute rate & -0.244 & -0.053 & -0.261 & -0.323 & +0.027 & -0.097 & +0.027 \\
Wikipedia token-frequency rank & +0.459 & +0.228 & +0.346 & +0.453 & +0.455 & +0.448 & +0.286 \\
Wikipedia TF/logit vocabulary agreement & +0.167 & -0.149 & -0.062 & -0.112 & -0.105 & +0.040 & +0.243 \\
\bottomrule
\end{tabular}}
\end{table}

\begin{table}[htbp]
\centering
\caption{Pearson $r$ with shared-split LRE faithfulness across 25 relations per model. Each relation has equal weight; larger lens scores mean better target ranks.}
\label{tab:shared-controls-pearson}
\scriptsize
\setlength{\tabcolsep}{3.5pt}
\resizebox{\linewidth}{!}{%
\begin{tabular}{lrrrrrrr}
\toprule
Predictor & OLMo & GPT-J & Mistral & Gemma & Qwen 8B & Qwen 27B & Pythia \\
\midrule
Model answer-token prior & +0.043 & +0.103 & -0.039 & +0.294 & +0.077 & +0.185 & -0.103 \\
Gold answer-token prior & +0.184 & +0.264 & +0.054 & +0.331 & +0.230 & +0.322 & +0.084 \\
Full-data samples/objects & +0.296 & +0.249 & +0.491 & +0.449 & +0.323 & +0.413 & +0.280 \\
Filtered samples/objects & +0.296 & +0.249 & +0.491 & +0.449 & +0.322 & +0.412 & +0.278 \\
Test samples/objects & +0.296 & +0.217 & +0.522 & +0.521 & +0.347 & +0.417 & +0.219 \\
Relation-prompt intercept-only agreement & +0.377 & +0.307 & +0.568 & +0.491 & +0.313 & +0.488 & +0.421 \\
Modal-target fraction (oracle) & +0.552 & +0.459 & +0.694 & +0.645 & +0.512 & +0.612 & +0.487 \\
Target entropy & -0.806 & -0.805 & -0.767 & -0.787 & -0.800 & -0.804 & -0.788 \\
Model gold-token accuracy & +0.683 & +0.721 & +0.505 & +0.548 & +0.603 & +0.475 & +0.660 \\
\bottomrule
\end{tabular}}
\end{table}

\begin{table}[htbp]
\centering
\caption{Pearson $r$ with shared-split LRE faithfulness across 25 relations per model. Document counts come from OLMo's published training split; they are an external corpus proxy for the other models.}
\label{tab:shared-corpus-pearson}
\scriptsize
\setlength{\tabcolsep}{3.5pt}
\resizebox{\linewidth}{!}{%
\begin{tabular}{lrrrrrrr}
\toprule
Predictor & OLMo & GPT-J & Mistral & Gemma & Qwen 8B & Qwen 27B & Pythia \\
\midrule
Subject document frequency & +0.409 & +0.461 & +0.035 & -0.004 & +0.358 & +0.342 & +0.510 \\
Object document frequency & +0.439 & +0.427 & +0.289 & +0.280 & +0.293 & +0.300 & +0.430 \\
Joint document frequency & +0.470 & +0.449 & +0.105 & +0.190 & +0.378 & +0.370 & +0.532 \\
Mean log(1+joint frequency) & +0.526 & +0.578 & +0.147 & +0.160 & +0.445 & +0.344 & +0.586 \\
Log(1+mean joint frequency) & +0.558 & +0.606 & +0.220 & +0.224 & +0.486 & +0.404 & +0.623 \\
P(object | subject) & +0.477 & +0.427 & +0.560 & +0.684 & +0.486 & +0.412 & +0.292 \\
P(subject | object) & -0.481 & -0.481 & -0.737 & -0.729 & -0.535 & -0.596 & -0.388 \\
Forward minus reverse conditional & +0.519 & +0.493 & +0.706 & +0.768 & +0.553 & +0.550 & +0.373 \\
Micro P(object | subject) & +0.560 & +0.490 & +0.497 & +0.611 & +0.557 & +0.322 & +0.437 \\
Micro P(subject | object) & +0.026 & -0.083 & -0.371 & -0.276 & -0.106 & -0.245 & -0.119 \\
\bottomrule
\end{tabular}}
\end{table}

\begin{table}[htbp]
\centering
\caption{Pearson $r$ with shared-split LRE faithfulness across 25 relations per model. Tuned rows select beta on the scored test rows and are exploratory.}
\label{tab:shared-beta-pearson}
\scriptsize
\setlength{\tabcolsep}{3.5pt}
\resizebox{\linewidth}{!}{%
\begin{tabular}{lrrrrrrr}
\toprule
Predictor & OLMo & GPT-J & Mistral & Gemma & Qwen 8B & Qwen 27B & Pythia \\
\midrule
Affine beta=0.5 & +0.421 & +0.116 & +0.408 & +0.485 & +0.319 & +0.353 & +0.156 \\
Affine beta=1 & +0.544 & +0.275 & +0.425 & +0.510 & +0.408 & +0.370 & +0.298 \\
Affine beta=1.5 & +0.632 & +0.412 & +0.439 & +0.541 & +0.485 & +0.385 & +0.393 \\
Affine beta=2 & +0.678 & +0.478 & +0.449 & +0.574 & +0.558 & +0.399 & +0.454 \\
Affine beta=3 & +0.697 & +0.514 & +0.458 & +0.622 & +0.620 & +0.419 & +0.530 \\
Affine beta=5 & +0.703 & +0.536 & +0.497 & +0.646 & +0.657 & +0.447 & +0.567 \\
Affine tuned top1 (in sample) & +0.404 & +0.276 & +0.401 & +0.485 & +0.406 & +0.345 & +0.306 \\
Affine tuned median-logrank (in sample) & +0.707 & +0.513 & +0.496 & +0.646 & +0.658 & +0.445 & +0.563 \\
\bottomrule
\end{tabular}}
\end{table}


\clearpage
\subsection{Complete relation-level measurements}
The following tables give every primary relation measurement. Exact values,
all beta settings, coverage, p-values for both correlation statistics,
paired intervals, and source hashes are retained in the accompanying CSV/JSON
artifacts under \texttt{code/results/shared\_baselines/}.
Historical unprefixed measurements, extended-beta convergence checks,
deduplicated Mistral grid-confirmation experiments, and causal/patching
analyses are catalogued separately; they are not mixed into the current
shared-split baseline tables.
\begin{table}[htbp]
\centering
\caption{OLMo 3 7B: all shared-split relation measurements. Lens columns are negative median log-ranks; Geva is mean attribute overlap in percent. D: differential Jacobian; J: raw Jacobian; LL: logit lens; DLL: differential logit lens; A: affine (beta=1); W: Wikipedia token-frequency rank.}
\label{tab:shared-relations-olmo3_7b}
\scriptsize
\setlength{\tabcolsep}{3.5pt}
\resizebox{\linewidth}{!}{%
\begin{tabular}{lrrrrrrrrr}
\toprule
Relation & $n$ & Faith & D & J & LL & DLL & A & Geva & W \\
\midrule
city\_in\_country & 6 & 0.833 & -0.540 & -0.540 & -3.395 & -3.168 & -3.300 & 10.662 & -1.161 \\
company\_ceo & 100 & 0.110 & -4.079 & -4.069 & -4.337 & -4.525 & -4.027 & 14.535 & -4.702 \\
company\_hq & 100 & 0.450 & -3.132 & -3.208 & -3.912 & -4.290 & -3.959 & 11.863 & -2.578 \\
country\_capital\_city & 8 & 1.000 & -0.878 & -0.724 & -3.497 & -3.128 & -3.445 & 12.952 & -1.791 \\
country\_currency & 14 & 0.357 & -3.208 & -2.971 & -3.566 & -3.288 & -3.767 & 9.439 & -4.707 \\
country\_language & 8 & 1.000 & -0.651 & -0.739 & -4.411 & -3.779 & -3.094 & 12.952 & -1.701 \\
country\_largest\_city & 8 & 1.000 & -0.301 & -0.151 & -3.716 & -3.495 & -3.669 & 13.382 & -2.756 \\
landmark\_in\_country & 100 & 0.800 & -1.892 & -1.993 & -3.988 & -4.561 & -3.579 & 8.675 & -1.836 \\
landmark\_on\_continent & 100 & 0.900 & -2.503 & -2.222 & -3.705 & -4.828 & -3.106 & 6.111 & -1.934 \\
name\_birthplace & 15 & 1.000 & -0.699 & -0.903 & -3.545 & -4.156 & -3.704 & 6.963 & -2.687 \\
name\_gender & 3 & 1.000 & -4.410 & -4.333 & -4.727 & -4.773 & -3.121 & 6.603 & -2.496 \\
name\_religion & 15 & 1.000 & -1.568 & -1.851 & -4.780 & -4.583 & -3.622 & 11.583 & -3.002 \\
person\_father & 100 & 0.060 & -3.740 & -3.531 & -4.419 & -4.547 & -3.925 & 12.123 & -2.479 \\
person\_mother & 100 & 0.120 & -4.045 & -4.116 & -4.469 & -4.500 & -3.793 & 12.712 & -3.000 \\
person\_native\_language & 100 & 0.930 & -1.567 & -0.954 & -4.354 & -4.306 & -3.108 & 5.168 & -1.522 \\
person\_occupation & 100 & 0.660 & -3.248 & -3.217 & -3.897 & -4.567 & -3.734 & 4.887 & -2.815 \\
person\_plays\_instrument & 100 & 0.740 & -2.264 & -2.102 & -4.288 & -4.613 & -2.624 & 5.649 & -2.235 \\
person\_plays\_position\_in\_sport & 100 & 0.860 & -2.712 & -2.747 & -4.316 & -4.312 & -4.067 & 3.141 & -1.914 \\
person\_plays\_pro\_sport & 100 & 0.910 & -0.977 & -1.146 & -3.729 & -4.011 & -2.551 & 7.672 & -2.135 \\
person\_university & 29 & 0.310 & -4.076 & -4.164 & -4.199 & -4.702 & -3.927 & 12.513 & -4.705 \\
pokemon\_evolutions & 20 & 0.250 & -3.191 & -3.135 & -4.519 & -4.812 & -3.582 & 4.558 & -2.129 \\
product\_by\_company & 100 & 0.690 & -1.361 & -1.097 & -4.152 & -4.323 & -3.260 & 10.348 & -1.402 \\
star\_constellation & 100 & 0.380 & -4.326 & -4.660 & -3.946 & -4.448 & -3.667 & 3.720 & -1.564 \\
superhero\_archnemesis & 55 & 0.345 & -3.611 & -3.492 & -4.363 & -4.772 & -3.718 & 13.294 & -2.149 \\
superhero\_person & 63 & 0.556 & -1.968 & -2.076 & -4.313 & -4.306 & -3.569 & 11.262 & -1.895 \\
\bottomrule
\end{tabular}}
\end{table}

\clearpage
\begin{table}[htbp]
\centering
\caption{GPT-J 6B: all shared-split relation measurements. Lens columns are negative median log-ranks; Geva is mean attribute overlap in percent. D: differential Jacobian; J: raw Jacobian; LL: logit lens; DLL: differential logit lens; A: affine (beta=1); W: Wikipedia token-frequency rank.}
\label{tab:shared-relations-gptj}
\scriptsize
\setlength{\tabcolsep}{3.5pt}
\resizebox{\linewidth}{!}{%
\begin{tabular}{lrrrrrrrrr}
\toprule
Relation & $n$ & Faith & D & J & LL & DLL & A & Geva & W \\
\midrule
city\_in\_country & 6 & 1.000 & -0.151 & -0.151 & -1.230 & -0.690 & -2.779 & 39.064 & -1.247 \\
company\_ceo & 100 & 0.090 & -3.569 & -3.587 & -4.122 & -4.070 & -3.903 & 35.761 & -4.401 \\
company\_hq & 100 & 0.390 & -2.221 & -2.322 & -3.135 & -3.420 & -3.813 & 24.930 & -4.401 \\
country\_capital\_city & 8 & 1.000 & -0.759 & -0.671 & -2.816 & -1.278 & -3.228 & 46.464 & -1.746 \\
country\_currency & 14 & 0.714 & -2.437 & -2.425 & -3.216 & -2.789 & -3.739 & 38.107 & -4.415 \\
country\_language & 8 & 0.875 & -0.841 & -0.778 & -3.016 & -2.052 & -2.196 & 46.464 & -1.674 \\
country\_largest\_city & 8 & 1.000 & -0.477 & -0.477 & -2.370 & -1.280 & -3.188 & 39.038 & -2.557 \\
landmark\_in\_country & 100 & 0.710 & -0.977 & -0.874 & -2.273 & -2.434 & -3.404 & 22.628 & -1.917 \\
landmark\_on\_continent & 100 & 0.890 & -1.877 & -1.880 & -3.308 & -4.449 & -3.238 & 15.712 & -1.915 \\
name\_birthplace & 15 & 0.867 & -1.204 & -1.342 & -2.947 & -2.838 & -3.715 & 33.773 & -2.734 \\
name\_gender & 3 & 1.000 & -4.319 & -4.393 & -3.952 & -4.529 & -3.491 & 20.841 & -2.502 \\
name\_religion & 15 & 0.867 & -1.944 & -1.732 & -3.730 & -3.796 & -3.766 & 32.355 & -3.002 \\
person\_father & 100 & 0.120 & -2.499 & -2.412 & -3.851 & -3.812 & -3.563 & 42.576 & -1.130 \\
person\_mother & 100 & 0.070 & -2.893 & -2.886 & -3.811 & -3.752 & -3.814 & 47.381 & -3.182 \\
person\_native\_language & 100 & 0.930 & -0.477 & -0.301 & -2.796 & -3.022 & -2.307 & 23.644 & -1.525 \\
person\_occupation & 100 & 0.650 & -3.869 & -3.980 & -4.074 & -4.068 & -4.174 & 16.758 & -4.403 \\
person\_plays\_instrument & 100 & 0.620 & -2.958 & -2.982 & -3.815 & -4.485 & -2.637 & 21.312 & -2.444 \\
person\_plays\_position\_in\_sport & 100 & 0.800 & -1.982 & -2.085 & -3.770 & -4.240 & -4.168 & 19.665 & -2.564 \\
person\_plays\_pro\_sport & 100 & 0.860 & -1.041 & -1.021 & -0.998 & -2.204 & -1.987 & 36.576 & -2.063 \\
person\_university & 25 & 0.520 & -3.913 & -3.817 & -3.228 & -3.702 & -3.829 & 48.173 & -2.310 \\
pokemon\_evolutions & 20 & 0.200 & -3.318 & -3.264 & -4.021 & -4.517 & -3.675 & 16.509 & -2.106 \\
product\_by\_company & 100 & 0.590 & -1.483 & -1.537 & -2.124 & -2.083 & -2.948 & 33.021 & -1.394 \\
star\_constellation & 100 & 0.550 & -3.323 & -2.781 & -3.698 & -4.363 & -3.735 & 8.931 & -1.728 \\
superhero\_archnemesis & 57 & 0.544 & -3.775 & -3.786 & -3.450 & -4.024 & -0.778 & 30.437 & -0.845 \\
superhero\_person & 63 & 0.492 & -1.415 & -1.431 & -3.999 & -4.332 & -3.434 & 30.047 & -1.892 \\
\bottomrule
\end{tabular}}
\end{table}

\clearpage
\begin{table}[htbp]
\centering
\caption{Mistral 7B: all shared-split relation measurements. Lens columns are negative median log-ranks; Geva is mean attribute overlap in percent. D: differential Jacobian; J: raw Jacobian; LL: logit lens; DLL: differential logit lens; A: affine (beta=1); W: Wikipedia token-frequency rank.}
\label{tab:shared-relations-mistral7b}
\scriptsize
\setlength{\tabcolsep}{3.5pt}
\resizebox{\linewidth}{!}{%
\begin{tabular}{lrrrrrrrrr}
\toprule
Relation & $n$ & Faith & D & J & LL & DLL & A & Geva & W \\
\midrule
city\_in\_country & 6 & 0.500 & -1.307 & -0.984 & -2.252 & -1.185 & -3.389 & 31.230 & -1.408 \\
company\_ceo & 100 & 0.040 & -3.940 & -3.873 & -4.106 & -3.981 & -3.753 & 36.625 & -4.217 \\
company\_hq & 100 & 0.560 & -2.097 & -2.155 & -2.492 & -2.693 & -3.759 & 28.815 & -2.168 \\
country\_capital\_city & 8 & 0.625 & -1.702 & -1.622 & -3.143 & -1.332 & -3.672 & 28.117 & -1.888 \\
country\_currency & 6 & 0.333 & -3.504 & -3.586 & -4.100 & -4.045 & -3.361 & 32.998 & -2.806 \\
country\_language & 8 & 0.625 & -0.151 & -0.151 & -1.073 & -0.540 & -3.195 & 28.117 & -1.816 \\
country\_largest\_city & 8 & 0.500 & -0.778 & -0.699 & -1.488 & -0.602 & -3.773 & 26.288 & -2.705 \\
landmark\_in\_country & 100 & 0.820 & -1.607 & -1.947 & -1.388 & -1.078 & -3.644 & 19.914 & -1.821 \\
landmark\_on\_continent & 100 & 0.920 & -2.860 & -2.842 & -3.911 & -3.917 & -2.961 & 12.602 & -1.605 \\
name\_birthplace & 15 & 1.000 & -0.699 & -1.146 & -1.978 & -2.041 & -3.817 & 16.130 & -2.760 \\
name\_gender & 3 & 0.667 & -3.397 & -2.842 & -4.385 & -4.439 & -2.956 & 8.986 & -2.196 \\
name\_religion & 15 & 1.000 & -3.117 & -2.866 & -4.140 & -3.985 & -3.619 & 15.524 & -3.026 \\
person\_father & 100 & 0.020 & -3.513 & -3.758 & -3.908 & -3.872 & -3.761 & 29.981 & -2.593 \\
person\_mother & 100 & 0.020 & -3.375 & -3.911 & -4.055 & -3.931 & -3.719 & 31.448 & -3.125 \\
person\_native\_language & 100 & 0.880 & -1.130 & -0.954 & -1.361 & -2.273 & -2.794 & 22.318 & -1.646 \\
person\_occupation & 100 & 0.790 & -2.175 & -1.819 & -2.994 & -2.953 & -3.392 & 16.757 & -2.567 \\
person\_plays\_instrument & 100 & 0.890 & -0.602 & -0.778 & -1.816 & -3.381 & -2.486 & 22.363 & -1.839 \\
person\_plays\_position\_in\_sport & 100 & 0.900 & -2.408 & -2.686 & -3.729 & -4.023 & -3.120 & 20.125 & -1.672 \\
person\_plays\_pro\_sport & 100 & 0.820 & -1.041 & -1.000 & -1.097 & -2.054 & -2.488 & 29.288 & -2.095 \\
person\_university & 23 & 0.304 & -3.885 & -3.889 & -3.551 & -3.783 & -3.872 & 34.317 & -3.172 \\
pokemon\_evolutions & 14 & 0.214 & -3.623 & -3.757 & -4.245 & -3.885 & -2.980 & 8.740 & -1.881 \\
product\_by\_company & 100 & 0.640 & -0.845 & -0.739 & -2.347 & -2.757 & -2.868 & 23.345 & -1.462 \\
star\_constellation & 100 & 0.270 & -3.815 & -3.837 & -4.135 & -4.109 & -3.309 & 6.859 & -1.105 \\
superhero\_archnemesis & 49 & 0.245 & -3.811 & -3.823 & -4.039 & -3.924 & -3.473 & 14.549 & -2.063 \\
superhero\_person & 59 & 0.203 & -3.222 & -3.119 & -4.136 & -4.044 & -3.674 & 10.866 & -2.143 \\
\bottomrule
\end{tabular}}
\end{table}

\clearpage
\begin{table}[htbp]
\centering
\caption{Gemma 2 9B: all shared-split relation measurements. Lens columns are negative median log-ranks; Geva is mean attribute overlap in percent. D: differential Jacobian; J: raw Jacobian; LL: logit lens; DLL: differential logit lens; A: affine (beta=1); W: Wikipedia token-frequency rank.}
\label{tab:shared-relations-gemma2_9b}
\scriptsize
\setlength{\tabcolsep}{3.5pt}
\resizebox{\linewidth}{!}{%
\begin{tabular}{lrrrrrrrrr}
\toprule
Relation & $n$ & Faith & D & J & LL & DLL & A & Geva & W \\
\midrule
city\_in\_country & 6 & 0.833 & -0.628 & -0.628 & -2.963 & -1.050 & -3.499 & 10.711 & -1.255 \\
company\_ceo & 100 & 0.060 & -4.782 & -4.704 & -4.736 & -4.932 & -4.129 & 24.595 & -5.108 \\
company\_hq & 100 & 0.480 & -2.713 & -2.689 & -4.152 & -4.223 & -3.939 & 17.790 & -2.134 \\
country\_capital\_city & 8 & 0.625 & -0.540 & -0.965 & -3.099 & -1.509 & -4.413 & 20.246 & -1.802 \\
country\_currency & 14 & 0.286 & -3.977 & -4.240 & -4.959 & -4.612 & -4.318 & 23.120 & -5.110 \\
country\_language & 8 & 1.000 & -1.784 & -2.177 & -1.922 & -1.751 & -3.804 & 20.246 & -1.747 \\
country\_largest\_city & 8 & 0.250 & -0.628 & -0.910 & -3.155 & -1.447 & -4.472 & 21.646 & -2.701 \\
landmark\_in\_country & 100 & 0.760 & -2.138 & -2.575 & -3.768 & -3.789 & -3.567 & 17.518 & -1.740 \\
landmark\_on\_continent & 100 & 0.930 & -2.593 & -2.223 & -4.498 & -5.114 & -2.733 & 9.698 & -1.964 \\
name\_birthplace & 15 & 1.000 & -1.708 & -1.643 & -2.461 & -4.742 & -3.912 & 13.690 & -2.662 \\
name\_gender & 3 & 0.667 & -4.658 & -4.563 & -3.722 & -4.531 & -3.227 & 17.503 & -2.584 \\
name\_religion & 15 & 0.933 & -3.082 & -2.623 & -4.774 & -4.462 & -4.006 & 21.134 & -5.111 \\
person\_father & 100 & 0.030 & -4.165 & -4.323 & -4.845 & -4.712 & -3.963 & 10.423 & -2.508 \\
person\_mother & 100 & 0.080 & -4.345 & -4.353 & -5.077 & -5.127 & -4.210 & 11.328 & -5.107 \\
person\_native\_language & 100 & 0.930 & -2.487 & -2.189 & -3.495 & -4.459 & -2.491 & 6.883 & -1.585 \\
person\_occupation & 100 & 0.810 & -3.155 & -2.996 & -3.264 & -4.483 & -3.541 & 5.856 & -2.380 \\
person\_plays\_instrument & 100 & 0.940 & -1.728 & -2.088 & -4.622 & -4.709 & -2.258 & 6.265 & -1.916 \\
person\_plays\_position\_in\_sport & 100 & 0.810 & -2.284 & -2.225 & -3.545 & -4.263 & -4.362 & 4.944 & -1.738 \\
person\_plays\_pro\_sport & 100 & 0.940 & -2.124 & -2.786 & -3.552 & -4.174 & -2.314 & 5.836 & -2.026 \\
person\_university & 31 & 0.226 & -3.989 & -4.233 & -4.720 & -4.373 & -4.217 & 11.807 & -2.994 \\
pokemon\_evolutions & 22 & 0.273 & -3.829 & -3.890 & -4.660 & -4.672 & -3.343 & 13.375 & -2.225 \\
product\_by\_company & 100 & 0.530 & -0.900 & -1.000 & -3.778 & -3.790 & -2.873 & 14.403 & -1.431 \\
star\_constellation & 100 & 0.750 & -4.019 & -3.685 & -4.928 & -3.533 & -3.522 & 1.468 & -1.923 \\
superhero\_archnemesis & 57 & 0.298 & -4.235 & -4.295 & -4.752 & -4.919 & -3.797 & 11.260 & -1.959 \\
superhero\_person & 63 & 0.476 & -3.109 & -2.866 & -4.368 & -4.346 & -3.483 & 11.984 & -1.924 \\
\bottomrule
\end{tabular}}
\end{table}

\clearpage
\begin{table}[htbp]
\centering
\caption{Qwen3 8B Base: all shared-split relation measurements. Lens columns are negative median log-ranks; Geva is mean attribute overlap in percent. D: differential Jacobian; J: raw Jacobian; LL: logit lens; DLL: differential logit lens; A: affine (beta=1); W: Wikipedia token-frequency rank.}
\label{tab:shared-relations-qwen3_8b_base}
\scriptsize
\setlength{\tabcolsep}{3.5pt}
\resizebox{\linewidth}{!}{%
\begin{tabular}{lrrrrrrrrr}
\toprule
Relation & $n$ & Faith & D & J & LL & DLL & A & Geva & W \\
\midrule
city\_in\_country & 6 & 0.833 & -0.452 & -0.452 & -0.000 & -0.151 & -3.406 & 24.275 & -1.399 \\
company\_ceo & 100 & 0.040 & -4.430 & -4.452 & -4.321 & -4.564 & -4.053 & 18.139 & -4.881 \\
company\_hq & 100 & 0.470 & -3.055 & -3.217 & -2.990 & -3.656 & -3.870 & 15.786 & -2.413 \\
country\_capital\_city & 8 & 1.000 & -1.050 & -1.146 & -2.866 & -0.651 & -3.856 & 19.596 & -1.807 \\
country\_currency & 14 & 0.357 & -3.662 & -3.994 & -4.038 & -3.955 & -3.970 & 21.572 & -4.885 \\
country\_language & 8 & 1.000 & -1.041 & -1.040 & -2.569 & -0.849 & -2.774 & 19.596 & -1.737 \\
country\_largest\_city & 8 & 0.875 & -0.845 & -0.952 & -2.304 & -0.874 & -3.868 & 16.972 & -2.742 \\
landmark\_in\_country & 100 & 0.720 & -2.182 & -2.308 & -2.462 & -2.622 & -3.567 & 15.602 & -1.772 \\
landmark\_on\_continent & 100 & 0.870 & -2.652 & -2.008 & -2.301 & -4.278 & -3.040 & 13.552 & -1.951 \\
name\_birthplace & 15 & 1.000 & -0.845 & -1.146 & -2.114 & -1.531 & -3.755 & 27.902 & -2.691 \\
name\_gender & 3 & 0.667 & -4.887 & -4.959 & -4.424 & -5.167 & -3.377 & 11.420 & -2.493 \\
name\_religion & 15 & 1.000 & -3.527 & -3.643 & -3.303 & -4.339 & -3.891 & 20.733 & -2.952 \\
person\_father & 100 & 0.040 & -4.142 & -4.324 & -4.141 & -4.386 & -3.815 & 26.589 & -2.966 \\
person\_mother & 100 & 0.060 & -4.348 & -4.418 & -4.062 & -4.531 & -3.839 & 25.081 & -3.205 \\
person\_native\_language & 100 & 0.920 & -2.057 & -2.023 & -2.582 & -4.061 & -2.771 & 16.061 & -1.597 \\
person\_occupation & 100 & 0.800 & -3.798 & -3.644 & -4.087 & -4.247 & -3.795 & 12.381 & -2.784 \\
person\_plays\_instrument & 100 & 0.720 & -3.483 & -2.987 & -3.295 & -4.832 & -2.526 & 16.039 & -2.071 \\
person\_plays\_position\_in\_sport & 100 & 0.800 & -1.675 & -1.618 & -3.424 & -4.402 & -3.954 & 14.014 & -1.911 \\
person\_plays\_pro\_sport & 100 & 0.960 & -1.826 & -1.712 & -1.518 & -3.856 & -2.586 & 23.160 & -2.052 \\
person\_university & 29 & 0.276 & -4.446 & -4.489 & -3.922 & -4.299 & -3.989 & 24.951 & -3.029 \\
pokemon\_evolutions & 20 & 0.150 & -4.434 & -4.242 & -4.187 & -4.672 & -3.569 & 8.432 & -2.143 \\
product\_by\_company & 100 & 0.660 & -1.743 & -1.527 & -2.884 & -3.013 & -2.888 & 19.515 & -1.431 \\
star\_constellation & 100 & 0.280 & -4.896 & -4.810 & -4.036 & -4.931 & -3.647 & 5.352 & -1.512 \\
superhero\_archnemesis & 55 & 0.382 & -4.104 & -3.867 & -4.178 & -4.695 & -3.533 & 21.269 & -2.271 \\
superhero\_person & 63 & 0.270 & -2.956 & -2.981 & -4.124 & -4.559 & -3.472 & 16.747 & -1.929 \\
\bottomrule
\end{tabular}}
\end{table}

\clearpage
\begin{table}[htbp]
\centering
\caption{Qwen3.8 27B: all shared-split relation measurements. Lens columns are negative median log-ranks; Geva is mean attribute overlap in percent. D: differential Jacobian; J: raw Jacobian; LL: logit lens; DLL: differential logit lens; A: affine (beta=1); W: Wikipedia token-frequency rank.}
\label{tab:shared-relations-qwen3_8_27b}
\scriptsize
\setlength{\tabcolsep}{3.5pt}
\resizebox{\linewidth}{!}{%
\begin{tabular}{lrrrrrrrrr}
\toprule
Relation & $n$ & Faith & D & J & LL & DLL & A & Geva & W \\
\midrule
city\_in\_country & 6 & 0.500 & -0.889 & -0.801 & -2.910 & -0.661 & -4.012 & 16.277 & -1.399 \\
company\_ceo & 100 & 0.010 & -4.715 & -4.690 & -4.435 & -4.831 & -4.196 & 11.426 & -5.095 \\
company\_hq & 100 & 0.490 & -2.486 & -2.506 & -4.190 & -4.657 & -4.026 & 8.295 & -2.665 \\
country\_capital\_city & 8 & 0.750 & -1.128 & -1.535 & -3.969 & -2.835 & -4.264 & 9.901 & -1.805 \\
country\_currency & 14 & 0.429 & -4.716 & -4.429 & -4.802 & -4.520 & -4.400 & 13.226 & -5.097 \\
country\_language & 8 & 1.000 & -2.185 & -2.468 & -3.566 & -3.048 & -3.320 & 9.901 & -1.734 \\
country\_largest\_city & 8 & 1.000 & -1.116 & -1.462 & -3.553 & -2.696 & -4.285 & 8.983 & -2.733 \\
landmark\_in\_country & 100 & 0.750 & -2.816 & -2.823 & -3.844 & -3.699 & -3.659 & 7.480 & -1.785 \\
landmark\_on\_continent & 100 & 0.960 & -2.969 & -2.049 & -4.570 & -5.147 & -3.197 & 6.735 & -1.916 \\
name\_birthplace & 15 & 0.867 & -1.041 & -1.114 & -3.379 & -3.534 & -3.840 & 17.031 & -2.882 \\
name\_gender & 3 & 0.667 & -4.979 & -4.973 & -4.630 & -5.274 & -3.219 & 9.480 & -2.485 \\
name\_religion & 15 & 1.000 & -2.770 & -2.759 & -4.707 & -4.597 & -3.932 & 15.501 & -2.945 \\
person\_father & 100 & 0.050 & -4.169 & -4.348 & -4.618 & -4.738 & -4.165 & 12.184 & -3.032 \\
person\_mother & 100 & 0.020 & -4.607 & -4.594 & -4.668 & -5.069 & -4.053 & 12.274 & -5.094 \\
person\_native\_language & 100 & 0.870 & -2.498 & -2.058 & -3.355 & -3.986 & -2.886 & 6.608 & -1.580 \\
person\_occupation & 100 & 0.810 & -4.044 & -3.718 & -4.677 & -5.081 & -3.652 & 5.660 & -2.741 \\
person\_plays\_instrument & 100 & 0.820 & -3.762 & -3.936 & -4.055 & -4.713 & -2.611 & 6.643 & -2.112 \\
person\_plays\_position\_in\_sport & 100 & 0.840 & -3.310 & -3.317 & -5.009 & -4.325 & -4.279 & 3.706 & -2.052 \\
person\_plays\_pro\_sport & 100 & 0.860 & -3.358 & -2.818 & -4.173 & -4.392 & -2.409 & 5.108 & -2.275 \\
person\_university & 31 & 0.419 & -4.703 & -4.385 & -4.202 & -5.036 & -4.014 & 10.167 & -3.023 \\
pokemon\_evolutions & 20 & 0.050 & -3.910 & -3.774 & -4.923 & -5.045 & -3.780 & 7.467 & -2.262 \\
product\_by\_company & 100 & 0.620 & -1.491 & -1.519 & -3.583 & -3.922 & -3.096 & 9.963 & -1.455 \\
star\_constellation & 100 & 0.270 & -4.517 & -4.462 & -4.837 & -5.228 & -3.556 & 0.986 & -1.740 \\
superhero\_archnemesis & 55 & 0.436 & -4.311 & -4.439 & -4.716 & -5.110 & -3.556 & 9.512 & -2.265 \\
superhero\_person & 63 & 0.143 & -3.322 & -3.344 & -4.466 & -4.675 & -3.796 & 8.866 & -1.927 \\
\bottomrule
\end{tabular}}
\end{table}

\clearpage
\begin{table}[htbp]
\centering
\caption{Pythia 6.9B: all shared-split relation measurements. Lens columns are negative median log-ranks; Geva is mean attribute overlap in percent. D: differential Jacobian; J: raw Jacobian; LL: logit lens; DLL: differential logit lens; A: affine (beta=1); W: Wikipedia token-frequency rank.}
\label{tab:shared-relations-pythia6_9b}
\scriptsize
\setlength{\tabcolsep}{3.5pt}
\resizebox{\linewidth}{!}{%
\begin{tabular}{lrrrrrrrrr}
\toprule
Relation & $n$ & Faith & D & J & LL & DLL & A & Geva & W \\
\midrule
city\_in\_country & 6 & 0.833 & -0.151 & -0.151 & -3.738 & -1.742 & -2.763 & 53.333 & -1.273 \\
company\_ceo & 100 & 0.090 & -3.998 & -4.113 & -3.380 & -4.035 & -3.736 & 66.887 & -4.402 \\
company\_hq & 100 & 0.320 & -4.057 & -4.287 & -3.332 & -3.570 & -3.891 & 42.868 & -2.636 \\
country\_capital\_city & 8 & 1.000 & -0.423 & -0.452 & -3.627 & -2.738 & -3.476 & 76.667 & -1.769 \\
country\_currency & 12 & 0.500 & -3.575 & -4.195 & -3.397 & -2.817 & -4.024 & 83.333 & -4.417 \\
country\_language & 8 & 1.000 & -0.690 & -0.845 & -3.403 & -3.222 & -2.738 & 76.667 & -1.677 \\
country\_largest\_city & 8 & 1.000 & -0.540 & -0.588 & -3.668 & -1.743 & -3.189 & 75.000 & -2.557 \\
landmark\_in\_country & 100 & 0.660 & -0.929 & -1.230 & -3.414 & -2.997 & -2.985 & 43.368 & -1.860 \\
landmark\_on\_continent & 100 & 0.910 & -2.928 & -3.435 & -4.245 & -4.345 & -2.992 & 32.898 & -1.918 \\
name\_birthplace & 15 & 0.867 & -0.477 & -0.477 & -3.807 & -3.512 & -3.667 & 58.333 & -2.709 \\
name\_gender & 3 & 1.000 & -4.123 & -4.436 & -2.265 & -4.295 & -3.284 & 50.000 & -2.523 \\
name\_religion & 15 & 0.933 & -1.146 & -1.623 & -4.366 & -4.514 & -3.574 & 56.061 & -3.008 \\
person\_father & 100 & 0.180 & -2.520 & -2.765 & -3.495 & -4.164 & -3.503 & 62.077 & -1.204 \\
person\_mother & 100 & 0.080 & -3.593 & -3.973 & -3.602 & -4.250 & -3.837 & 69.048 & -4.402 \\
person\_native\_language & 100 & 0.820 & -0.977 & -1.241 & -3.024 & -3.682 & -2.229 & 43.841 & -1.568 \\
person\_occupation & 100 & 0.750 & -4.547 & -4.587 & -4.177 & -4.238 & -4.235 & 30.468 & -4.404 \\
person\_plays\_instrument & 100 & 0.600 & -2.444 & -2.614 & -4.091 & -4.314 & -2.895 & 38.782 & -2.941 \\
person\_plays\_position\_in\_sport & 100 & 0.840 & -4.199 & -4.228 & -4.096 & -4.446 & -4.093 & 45.401 & -1.911 \\
person\_plays\_pro\_sport & 100 & 0.790 & -0.812 & -1.000 & -3.641 & -3.919 & -2.130 & 54.630 & -2.051 \\
person\_university & 25 & 0.440 & -3.942 & -4.242 & -2.704 & -4.259 & -3.623 & 44.118 & -2.451 \\
pokemon\_evolutions & 16 & 0.125 & -3.899 & -4.538 & -3.056 & -4.272 & -3.344 & 34.524 & -3.411 \\
product\_by\_company & 100 & 0.470 & -2.553 & -3.351 & -3.738 & -4.001 & -2.987 & 64.289 & -1.398 \\
star\_constellation & 100 & 0.150 & -3.880 & -3.982 & -3.244 & -4.320 & -3.193 & 47.422 & -1.459 \\
superhero\_archnemesis & 55 & 0.455 & -3.411 & -3.547 & -3.107 & -4.136 & -3.258 & 61.689 & -1.756 \\
superhero\_person & 61 & 0.443 & -2.717 & -3.419 & -3.848 & -4.322 & -3.483 & 67.100 & -1.922 \\
\bottomrule
\end{tabular}}
\end{table}

\clearpage
